\phantomsection
\addcontentsline{toc}{subsection}{Inciso a)}
\subsection*{Inciso a)}

% Seccion con nombres de clase largos, que no se parten con guion.
\sloppy

\textit{Identifica el patrón de diseño más adecuado, justifica tu selección basándote en las
características del problema y explica la propuesta de solución considerando las clases que
intervienen y las relaciones entre ellas. ¿Por qué no usar Mediator o un simple método
estático?}

\vspace{0.4cm}

\subsubsection*{Lectura del enunciado}

Cuatro preguntas antes de abrir el catálogo, porque la decisión sale de ellas:

\begin{itemize}[leftmargin=1.2cm]
    \item \textbf{¿Qué objetos existen?} Seis componentes independientes
          (\texttt{Amplificador}, \texttt{ReproductorDvd}, \texttt{ConsolaJuegos},
          \texttt{Proyector}, \texttt{Pantalla}, \texttt{Luces}) y un cliente que quiere
          hacer tres cosas.
    \item \textbf{¿Dónde está la dificultad?} No en los componentes, que funcionan. En que
          una actividad corriente exige once llamadas a seis objetos distintos
          \emph{en un orden que nadie ha escrito}.
    \item \textbf{¿Qué pide el enunciado que no cambie?} Nada lo dice con esas palabras,
          pero los componentes se describen como ``complejos'' y preexistentes: la solución
          no debería obligar a reescribirlos.
    \item \textbf{¿Qué se anuncia para después?} Actividades nuevas sin tocar la solución
          (inciso c), notificaciones desde los componentes (inciso d) y secuencias definidas
          en un archivo (restricción).
\end{itemize}

La segunda respuesta es la que decide. El problema no es de \emph{comportamiento} ni de
\emph{creación}: es que el conocimiento sobre cómo se usa el subsistema está repartido por
todos sus clientes.

\subsubsection*{Patrón seleccionado}

El patrón adecuado es \textbf{Facade} (Fachada), patrón estructural del catálogo
\textit{Gang of Four}. Su intención declarada es: \textit{proporcionar una interfaz unificada
para un conjunto de interfaces de un subsistema, definiendo una interfaz de más alto nivel que
hace al subsistema más fácil de usar}. Esa frase es el enunciado con otras palabras.

\subsubsection*{Justificación a partir de las características del problema}

Tres rasgos del problema apuntan al mismo sitio:

\begin{enumerate}[leftmargin=1.2cm]
    \item \textbf{La complejidad es de uso, no de implementación.} Ningún componente es
          difícil por sí mismo: \texttt{Pantalla} baja y sube. Lo difícil es la
          \emph{coreografía}, y una coreografía no tiene dónde vivir si no se le da una
          clase.
    \item \textbf{La dependencia tiene una sola dirección.} El cliente necesita al
          subsistema; el subsistema no necesita al cliente ni se necesita a sí mismo: la
          pantalla no consulta al amplificador. Esa asimetría es exactamente la que Facade
          supone, y es lo que lo separa de Mediator.
    \item \textbf{Hay un vocabulario del dominio que el código no nombra.} ``Ver una
          película'' existe para el usuario y no existe en el sistema. La fachada no es una
          utilidad técnica: es el sitio donde ese vocabulario se vuelve código.
\end{enumerate}

Conviene añadir el criterio de diseño que respalda la elección: el
\textbf{principio de mínimo conocimiento}, o ley de Deméter. Antes, el cliente conversa con
seis objetos y conoce el orden entre ellos; después, conversa con uno. El acoplamiento no
desaparece, se \emph{concentra}: queda un único punto que conoce el subsistema, y ese punto
es el que hay que tocar cuando el subsistema cambie.

\subsubsection*{La solución ingenua, y qué rompe}

La solución sin patrón es el cliente haciendo el trabajo:

{\small
\begin{verbatim}
// En cada sitio desde donde alguien quiera ver una pelicula:
luces.Atenuar(10);
pantalla.Bajar();
proyector.Encender();
proyector.Entrada("dvd");
proyector.Formato("16:9");
amplificador.Encender();
amplificador.Entrada("dvd");
amplificador.ModoEnvolvente(true);
amplificador.PonerVolumen(7);
dvd.Encender();
dvd.Reproducir("Blade Runner");
\end{verbatim}
}

Funciona, y por eso es tentadora. Lo que rompe:

\begin{enumerate}[leftmargin=1.2cm]
    \item \textbf{El conocimiento se duplica.} La secuencia se repite en cada cliente: el
          mando, la aplicación del móvil, la rutina de la noche. Once líneas por tres
          clientes y tres actividades son noventa y nueve ocasiones de equivocarse.
    \item \textbf{El orden no está escrito en ningún sitio.} Encender el proyector antes de
          bajar la pantalla no da un error de compilación: da una proyección sobre la pared.
          Es un requisito real que vive solo en la cabeza de quien escribió el cliente.
    \item \textbf{Un cambio en el subsistema alcanza a todos los clientes.} Si el
          amplificador nuevo exige fijar la entrada antes de encenderse, hay que buscar
          todas las copias.
    \item \textbf{El cliente depende de seis clases} en lugar de una, y de los detalles de
          cada una: los nombres de las entradas, los rangos de volumen, los formatos.
    \item \textbf{No hay dónde colgar lo que viene después.} Sin un objeto que represente la
          actividad, los incisos (c) y (d) y la restricción no tienen sujeto: no se puede
          configurar desde un archivo una secuencia que no existe como dato.
\end{enumerate}

Y la nota justa: con un solo cliente y una sola actividad, estas once líneas son la solución
correcta y cualquier patrón sería sobreingeniería. Lo que las descalifica aquí es lo que el
enunciado añade a continuación.

\subsubsection*{Qué patrón resuelve esto}

El problema es ``dar una interfaz simple a un conjunto de objetos que ya funcionan, sin
tocarlos''. Los candidatos razonables, con su veredicto:

\begin{table}[h!]
    \centering
    \small
    \begin{tabular}{p{2.9cm}p{12.1cm}}
        \toprule
        \textbf{Patrón} & \textbf{¿Sirve para este problema?} \\
        \midrule
        Adapter & \textbf{No:} convierte una interfaz en \emph{otra} que el cliente ya
                  espera, normalmente de una sola clase. No simplifica un conjunto ni
                  inventa operaciones nuevas. Se confunde con Facade porque los dos
                  envuelven, pero la intención es distinta: adaptar no es simplificar. \\
        \addlinespace
        Proxy & \textbf{No:} mantiene \emph{la misma} interfaz y controla el acceso a ella.
                Aquí la interfaz nueva es justamente el punto. \\
        \addlinespace
        Mediator & \textbf{No:} coordina la colaboración entre objetos que se comunican
                   \emph{entre sí}. Aquí no se comunican. Se detalla más abajo. \\
        \addlinespace
        Builder & \textbf{No:} construye paso a paso un objeto complejo. Suena por lo de
                  ``paso a paso'', pero lo que se construye aquí no es un objeto: es una
                  secuencia de órdenes sobre objetos que ya existen. \\
        \addlinespace
        Command & \textbf{Parcial:} convierte cada actividad en un objeto, que es lo que
                  hace falta para añadir actividades sin tocar la fachada. Es la respuesta
                  del inciso (c), y se combina con el elegido en lugar de sustituirlo. \\
        \addlinespace
        Template Method & \textbf{No:} fija el esqueleto de un algoritmo y deja los pasos a
                          las subclases. Las tres actividades no comparten esqueleto: usan
                          componentes distintos en órdenes distintos. \\
        \addlinespace
        \textbf{Facade} & \textbf{Sí:} define una interfaz de alto nivel sobre el
                          subsistema, con un método por actividad del dominio, sin modificar
                          ningún componente. \\
        \bottomrule
    \end{tabular}
\end{table}

\subsubsection*{Por qué no Mediator}

Los dos patrones ponen un objeto en el medio, de ahí la confusión. Se diferencian en cuatro
puntos decisivos:

\begin{enumerate}[leftmargin=1.2cm]
    \item \textbf{Topología.} Mediator resuelve una maraña $N \leftrightarrow N$: colegas
          que se conocen entre sí y a los que se desacopla haciéndolos hablar a través de un
          intermediario. Facade resuelve un acceso $1 \rightarrow N$: un cliente externo
          frente a un subsistema cuyos miembros \emph{no} se hablan. En este problema la
          pantalla no consulta al amplificador, ni las luces al proyector. La maraña que
          Mediator deshace aquí no existe.
    \item \textbf{Quién se modifica.} En Mediator cada colega guarda una referencia al
          mediador y le notifica sus eventos, de modo que \textbf{hay que modificar los seis
          componentes}. Facade no toca ninguno. Dado que el enunciado los describe como
          componentes complejos que ya existen, esa diferencia no es de estilo.
    \item \textbf{Dirección del flujo.} En Facade las órdenes van hacia dentro y el
          subsistema no responde. Cuando sí se quiera que un componente avise de algo, el
          patrón correcto no es Mediator sino \textbf{Observer}, que es exactamente lo que
          pide el inciso (d): notificar cambios de estado sin que el componente sepa a quién.
    \item \textbf{Qué le queda al cliente.} Un mediador que se limita a reenviar mensajes no
          le ahorra nada: el cliente seguiría pidiendo la película paso a paso. Lo que el
          problema necesita es una operación nueva de alto nivel, y eso es Facade.
\end{enumerate}

Dicho esto, la frontera es real: una fachada que empezara a recibir eventos de los
componentes y a reaccionar coordinándolos \emph{se estaría convirtiendo} en un mediador. El
inciso (d) roza ese límite y conviene no cruzarlo sin decirlo.

\subsubsection*{Por qué no un simple método estático}

Un \texttt{static void VerPelicula(...)} haría lo mismo en la primera versión. Deja de hacerlo
en cuanto el enunciado sigue:

\begin{itemize}[leftmargin=1.2cm]
    \item \textbf{No tiene estado.} La fachada guarda las referencias a los seis
          componentes. El método estático tendría que recibirlos como seis parámetros en
          cada llamada, devolviéndole al cliente justo el conocimiento que se le quería
          quitar, o leerlos de variables globales.
    \item \textbf{No es sustituible.} No se puede inyectar una fachada de mentira en una
          prueba, ni montar dos salas con componentes distintos. Un método estático no se
          implementa: se llama.
    \item \textbf{No es polimórfico.} No admite una variante ``sala premium'' ni un
          decorador que registre lo que se hace. Con un objeto, las dos cosas son naturales.
    \item \textbf{No se puede configurar.} La restricción final pide definir las secuencias
          en un archivo YAML. Un método estático lleva la secuencia escrita en su cuerpo;
          un objeto puede recibirla al construirse. La restricción, por sí sola, descarta la
          opción.
    \item \textbf{No tiene ciclo de vida.} El inciso (d) pide que la fachada sea un
          Singleton; un método estático no se instancia, así que no hay instancia que
          controlar, ni forma de inicializarla con lo que haga falta.
\end{itemize}

Otra vez, en honor a la verdad: para un cine único, con tres actividades fijas y sin pruebas,
el método estático bastaría. Lo que lo descalifica no es el purismo, son los incisos
siguientes.

\subsubsection*{Propuesta de solución: clases y relaciones}

\textbf{El subsistema} (seis componentes, no se modifican):

\begin{itemize}[leftmargin=1.2cm]
    \item \texttt{IComponente}: contrato mínimo con \texttt{Nombre} y \texttt{Estado}, que
          sirve para la traza y para el Observer del inciso (d).
    \item \texttt{ComponenteBase}: clase base con el estado y el registro de lo que ocurre.
    \item \texttt{Amplificador}, \texttt{ReproductorDvd}, \texttt{ConsolaJuegos},
          \texttt{Proyector}, \texttt{Pantalla} y \texttt{Luces}, cada uno con sus
          operaciones propias y tipadas.
\end{itemize}

\textbf{La fachada} (una clase, el único punto que conoce la coreografía):

\begin{itemize}[leftmargin=1.2cm]
    \item \texttt{HomeTheaterFacade}: recibe los seis componentes al construirse y publica
          un método por actividad: \texttt{VerPelicula}, \texttt{TerminarPelicula},
          \texttt{EscucharMusica} y \texttt{JugarVideojuego}.
\end{itemize}

\begin{figure}[h!]
    \centering
    \begin{tikzpicture}[
        font=\small,
        caja/.style={draw, rounded corners=2pt, align=center, inner sep=4pt},
        iface/.style={caja, fill=black!5},
        usa/.style={draw, -{Stealth[length=2mm]}},
        implementa/.style={draw, dashed, -{Triangle[open, length=3mm, width=2.5mm]}},
        nota/.style={font=\footnotesize\itshape}
    ]

    \node[caja] (cliente) at (0,0) {Cliente\\\footnotesize (mando, app, rutina)};
    \node[caja] (facade) at (5.2,0)
        {\texttt{HomeTheaterFacade}\\[2pt]
         \footnotesize\texttt{VerPelicula(titulo)}\\[-1pt]
         \footnotesize\texttt{EscucharMusica(lista)}\\[-1pt]
         \footnotesize\texttt{JugarVideojuego(juego)}};

    \draw[usa] (cliente.east) -- node[midway, above, nota] {1 llamada} (facade.west);

    \node[caja] (c1) at (10.8,1.9) {\texttt{Amplificador}};
    \node[caja] (c2) at (10.8,1.0) {\texttt{ReproductorDvd}};
    \node[caja] (c3) at (10.8,0.1) {\texttt{ConsolaJuegos}};
    \node[caja] (c4) at (10.8,-0.8) {\texttt{Proyector}};
    \node[caja] (c5) at (10.8,-1.7) {\texttt{Pantalla}};
    \node[caja] (c6) at (10.8,-2.6) {\texttt{Luces}};

    \foreach \c in {c1,c2,c3,c4,c5,c6} {
        \draw[usa] (facade.east) -- (\c.west);
    }

    \node[iface] (icomp) at (10.8,3.1) {\textit{«interface»}\\\texttt{IComponente}};
    \draw[implementa] (c1.north) -- (icomp.south);

    \node[nota, align=center] at (5.2,-2.4)
        {las seis flechas van en un solo sentido:\\el subsistema no conoce a la fachada};

    \end{tikzpicture}
    \caption{La fachada concentra el conocimiento del orden. La dependencia va en una sola
    dirección, y por eso ningún componente se modifica.}
\end{figure}

\textbf{Relaciones.} La fachada \emph{conoce} a los seis componentes y los seis
\emph{ignoran} que existe. No hay herencia entre ellos ni entre la fachada y nada: Facade es
un patrón de composición pura. El cliente depende de la fachada, y puede además seguir
dependiendo de un componente concreto si lo necesita.

\subsubsection*{Lo que el patrón no hace, dicho antes de que se note}

\begin{itemize}[leftmargin=1.2cm]
    \item \textbf{No encapsula el subsistema.} Los componentes siguen siendo públicos y
          accesibles. Es deliberado y está en el catálogo: la fachada es un atajo, no una
          aduana. Quien necesite algo que ella no ofrece, baja un nivel. La parte 5 de la
          demostración lo enseña.
    \item \textbf{No reduce el acoplamiento total, lo concentra.} Si el amplificador cambia
          su interfaz, hay exactamente un archivo que corregir; pero hay que corregirlo.
    \item \textbf{Puede degenerar.} Una fachada que acumula métodos hasta volverse el único
          sitio donde pasan cosas es un objeto dios. El inciso (c) ataca justamente esa
          deriva.
    \item \textbf{No añade comportamiento.} Cada paso ya existía. Lo que aporta es el orden,
          que es conocimiento, no código nuevo.
\end{itemize}

\subsubsection*{Observación sobre la combinación con otros patrones}

Los incisos siguientes no cambian de patrón: lo completan. \textbf{Command} convierte cada
actividad en un objeto, de modo que añadir una no obligue a modificar la fachada (inciso c).
\textbf{Singleton} controla que haya una sola fachada por sala y \textbf{Observer} permite que
los componentes avisen de sus cambios de estado sin saber quién escucha (inciso d). La
restricción final cierra el círculo: si una actividad es un objeto construible, su secuencia
puede venir de un archivo en lugar de estar escrita en el código.

\fussy
